\section{ARIS-Rotate Dataset}
\label{sec:dataset}

\subsection{Construction}
\label{sec:dataset:construction}

We apply the \aris{} pipeline to construct \textbf{ARIS-Rotate},
a dataset for rotation-conditioned image editing.
Each sample in ARIS-Rotate pairs a canonical front-view image of an object
with a target-view image and a natural-language editing instruction.

\paragraph{Object selection and diversity.}
We generate 50 objects spanning diverse categories including furniture, household items,
tools, decorative objects, and outdoor equipment.
Object descriptions are generated to cover a range of materials (wood, metal, plastic, fabric),
shapes (compact, elongated, symmetric, asymmetric),
and color profiles.
This diversity is intended to stress-test the rendering pipeline
across a range of 3D reconstruction and lighting scenarios.

\paragraph{Viewpoint specification.}
Each object is rendered at eight equally-spaced horizontal azimuths.
We represent viewpoints using natural-language labels rather than numeric angles,
reflecting how users would naturally describe rotation instructions:

\begin{center}
\begin{tabular}{cl}
\toprule
Azimuth & View Label \\
\midrule
$0^\circ$   & front view \\
$45^\circ$  & front-right view \\
$90^\circ$  & right side view \\
$135^\circ$ & back-right view \\
$180^\circ$ & back view \\
$225^\circ$ & back-left view \\
$270^\circ$ & left side view \\
$315^\circ$ & front-left view \\
\bottomrule
\end{tabular}
\end{center}

\paragraph{Editing pairs and instruction templates.}
We designate the $0^\circ$ (front view) render as the canonical source image.
For each of the remaining seven target views,
we form an editing pair with the following instruction template:

\begin{center}
\texttt{Rotate this object from front view to \{target\_view\}.}
\end{center}

Each of the 50 objects contributes 7 editing pairs,
for a total of \textbf{350 training pairs}.
The full rendered image set comprises \textbf{400 images} (50 objects $\times$ 8 views).

\begin{figure}[t]
  \centering
  % TODO: Insert dataset sample grid figure
  \fbox{\parbox{0.95\linewidth}{\centering\vspace{3cm}
    \textbf{[Figure 3: Dataset Sample Grid]}\\
    Grid: rows = 4 representative objects,
    columns = 8 viewpoints (front view $\to$ front-left view).\\
    Each cell shows the rendered image at that viewpoint after VLM evolution.
    \vspace{3cm}}}
  \caption{
    Sample images from ARIS-Rotate.
    Each row shows one object rendered at all eight horizontal viewpoints,
    from front view ($0^\circ$) to front-left view ($315^\circ$).
    All images have been approved by the VLM reviewer through the evolution loop.
  }
  \label{fig:dataset_grid}
\end{figure}

\subsection{Data Split}
\label{sec:dataset:split}

We partition objects using an \emph{object-disjoint} split:
the same object cannot appear in both training and test sets.
This ensures that evaluation measures generalization to unseen objects,
not memorization of seen objects from different viewpoints.

\begin{center}
\begin{tabular}{lccc}
\toprule
Split & Objects & Pairs & Images \\
\midrule
Train & 35 & 245 & 280 \\
Val   & 5  & 35  & 40  \\
Test  & 10 & 70  & 80  \\
\midrule
\textbf{Total} & \textbf{50} & \textbf{350} & \textbf{400} \\
\bottomrule
\end{tabular}
\end{center}

\subsection{Dataset Statistics}
\label{sec:dataset:stats}

\tbd{Fill in after data generation completes.}
Statistics to report:
\begin{itemize}
  \item Distribution of VLM evolution rounds per object (histogram)
  \item Distribution of final VLM quality scores
  \item Fraction of objects requiring $\geq$1, $\geq$3, $\geq$5 refinement rounds
  \item Breakdown of primary issue types diagnosed across all objects
    (flat\_lighting, color\_shift, scene\_light\_mismatch, etc.)
\end{itemize}
